\documentclass[10pt,a4paper]{amsart}
\usepackage[margin=19mm]{geometry}
\usepackage{amsmath,amssymb,mathtools}
\usepackage[scr=rsfs,cal=euler]{mathalfa}
\usepackage{xfrac}
\usepackage[unicode,hidelinks,bookmarks=false]{hyperref}
\newcommand{\ie}{i.e.\ }
\newcommand{\cf}{cf.\ }
\newcommand{\resp}{respectively\ }
\newcommand{\Subsection}{\subsection}
\providecommand{\nobreakdash}{\nobreak\hbox{-}}
\setlength{\emergencystretch}{2em}
\multlinegap0pt
\usepackage{csquotes}
\newtheorem{theorem}{Theorem}
\newtheorem{corollary}{Corollary}
\DeclareMathOperator{\esssup}{esssup}
\title{On equality of the $L^\infty$ norm of the gradient under the Hausdorff and Lebesgue measure}
\author{Ng Ze-An}
\date{Source: arXiv:2505.08010v2 (CC BY 4.0)}
\begin{document}
\maketitle

\noindent\emph{Source and licence.} On equality of the $L^\infty$ norm of the gradient under the Hausdorff and Lebesgue measure. Version arXiv:2505.08010v2 (CC BY 4.0), distributed by arXiv under the Creative Commons Attribution 4.0 International licence, http://creativecommons.org/licenses/by/4.0/, as recorded in the arXiv licence metadata for this exact version and shown on the abstract page https://arxiv.org/abs/2505.08010v2. Full licence notice: \texttt{licenses/arxiv-2505.08010-CC-BY-4.0.txt}.\par\smallskip

\noindent\textit{Editorial scope.} Counted unit: the article's theorem maintheorem (source0.tex lines 43--50). The article's complete original proof of it is delivered with it (source0.tex lines 341--383, one \texttt{proof} environment of 43 lines, which the article places away from the statement). No mathematics is written, repaired or re-derived: the statement and the proof are the article's own lines, in the article's own notation, and no retained line is substituted or re-flowed.  Retained uncounted as the source-local context the proof needs: lines 59--61.  Nothing was omitted from any retained span, so the package carries no omission record.  No retained line contains a citation, so the wrapper carries no bibliography and no bibliography entry.  No retained line contains a figure environment, an image-inclusion command or drawing code, so no image is delivered and the package declares no asset dependency.  Editorial formatting only, in addition: an AMS \texttt{amsart} preamble that restates the article's own declarations and macros used by the retained lines, namely the environments \texttt{theorem}, \texttt{corollary} on separate counters, together with the standard packages those lines need.  The retained environments print in this document as Theorem 1, Corollary 1 in order of appearance, whereas the article prints the same items with the numbering of its own full document.  The complete native source of the article as distributed in the arXiv e-print for this exact version is bundled unchanged as \texttt{sources/178001/source0.tex}.
\begin{corollary}\label{corollary2}
Let $f_n: \Omega \to \mathbb R$ be everywhere differentiable. Assume that $f_n \to f$ in $W^{1, \infty}$. Then $f$ is everywhere differentiable and further $f_n \to f$ and $\nabla f_n \to \nabla f$ uniformly.
\end{corollary}

\begin{theorem}\label{maintheorem}
Let $\Omega$ be an open subset of $\mathbb R^n$, and let $f: \Omega \to \mathbb R$ be continuous, and further differentiable $\mathcal H^s$-almost everywhere, for some real $0 \leq s \leq n-1$. Then we have

\begin{equation}\label{eq1}\sup_{\Omega} |\nabla f| = \underset{\Omega}{\esssup} \, |\nabla f|\end{equation}

\noindent where the essential supremum is taken with respect to Lebesgue measure. In particular,
\begin{equation}\label{eq2}\|\nabla f\|_{L^\infty (\mathcal H^s)} = \|\nabla f\|_{L^\infty(\mathcal H^n)}.\end{equation}
\end{theorem}

\begin{proof}[Proof of Proposition \ref{corollary2}] 

Since $f_n \to f$ in $W^{1, \infty}$, we have that 
$$\limsup_{n, m \to \infty} \, \underset{\Omega}{\esssup} \, |\nabla f_n - \nabla f_m| \to 0,$$
and thus, by Theorem \ref{maintheorem}, also
$$\limsup_{n, m \to \infty} \,  \sup_{\Omega} |\nabla f_n - \nabla f_m| \to 0.$$

\noindent That is, the sequence $\{\nabla f_n\}$ is Cauchy in supremum norm and hence converges uniformly to some $g: \Omega \to \mathbb R^n$. We claim now that $f$ is differentiable with $\nabla f = g$.

First, we show that the functions $f_n$ are uniformly differentiable - that is, for each $x \in \mathbb R^n$, we have that for every $\varepsilon > 0$ there exists $\delta > 0$ such that for all $n \in \mathbb N$,
\begin{equation}\label{unifdif}\frac{|f_n (x+y) - f_n (x) - \nabla f_n(x) \cdot y|}{y} < \varepsilon\end{equation}
whenever $|y| < \delta$. To this end, let $\varepsilon > 0$ be arbitrary. We pick $N \in \mathbb N$ so large such that
$$\sup_{n, m \geq N} \sup_{\Omega} |\nabla f_n - \nabla f_m| \leq \varepsilon.$$
Then for all $n \geq N$, we have that
$$\begin{aligned}&|f_n (x+y) - f_n (x) - \nabla f_n(x) \cdot y| \\
& \leq |f_n (x+y) - f_n (x) - (f_N (x+y) - f_N (x))| + |f_N (x+y) - f_N (x) - \nabla f_N (x) \cdot y| \\
&+ |(\nabla f_N(x) - \nabla f_n (x)) \cdot y| \\
& \leq \sup_\Omega |\nabla(f_n - f_N)(z)| \,  |y| + |f_N (x+y) - f_N (x) - \nabla f_N (x) \cdot y| + \varepsilon |y| \\
& \leq 2 \varepsilon |y| + |f_N (x+y) - f_N (x) - \nabla f_N (x) \cdot y|,\end{aligned}$$
so
$$\frac{|f_n (x+y) - f_n (x) - \nabla f_n (x)\cdot y|}{|y|} \leq 2\varepsilon + \frac{|f_N (x+y) - f_N (x) - \nabla f_N (x)|}{|y|}.$$
By differentiability of $f_N$, there is some $\delta_N$ such that the last term abovw is less than $\varepsilon$ for all $|y| < \delta_N$. Similarly, for each $0 \leq i < N$, there exists some $\delta_i$ such that 
$$\frac{|f_i (x+y) - f_i (x) - \nabla f_i (x) \cdot y|}{|y|} < \varepsilon$$
for all $|y| < \delta_i.$ Since $\varepsilon > 0$ was arbitrary, setting $\delta := \min_{0 \leq i \leq N} \delta_i$, we conclude.


Now we show the differentiability of $f$. Let $x \in \Omega$, $\varepsilon > 0$ be arbitrary. By the earlier discussion, there exists some $\delta > 0$ such that

$$\frac{|f_n (x+y) - f_n (x) - \nabla f_n (x) \cdot y|}{|y|} \leq \varepsilon$$
for all $y < \delta$.

Fix a small parameter $0 < h < 1$ which we will later send to $0$. Fix $n \in \mathbb N$ so large such that $\sup_{\Omega} |\nabla f_n - \nabla f| < \varepsilon$ and $\sup_{\Omega} |f_n - f| < h \delta \varepsilon$. Then for all $y \in B_{\delta}(x) \setminus B_{h\delta}(x)$, we compute using (\ref{unifdif}):
$$\begin{aligned}&|f(x+y) - f(x) - g(x) \cdot y|\\&\leq |f(x+y) - f_n (x+y)| + |f(x) - f_n (x)| + |f_n (x+y) - f_n (x) - \nabla f_n (x) \cdot y| \\
&+ |(\nabla f_n (x) - g) \cdot y|,\end{aligned} $$
so
$$\begin{aligned}\frac{|f(x+y) - f(x) - g(x) \cdot y|}{|y|} & \leq 2 \varepsilon + \varepsilon + \varepsilon \\
& = 4\varepsilon.\end{aligned}$$
We reiterate that the above holds for all $y \in B_{\delta} \setminus B_{h\delta}$.

Now fix a sequence of positive numbers $h_i \to 0_+$. Repeating the above argument with $h_i$ in place of $h$, and consequently $n_i$ in place of $n$, we obtain that for all $i$, and for all $y \in B_\delta (x) \setminus B_{h_i \delta}(x)$, we have
$$\frac{|f(x+y) - f(x) - g(x) \cdot y|}{|y|} < 4\varepsilon.$$
Since $h_i \to 0$, the above actually holds for all $y \in B_\delta(x) \setminus \{0\}$, which, since $\varepsilon > 0$ was arbitrary, shows the differentiability of $f$ at $x$ with derivative $g(x)$, as was to be shown.
\end{proof}

\end{document}
